\documentclass[11pt]{amsart}

\usepackage[T1]{fontenc}
\usepackage{lmodern}
\usepackage{microtype}
\usepackage{amsmath,amssymb,amsthm}
\usepackage[hidelinks]{hyperref}

\hypersetup{
  pdftitle={Counterexamples to Zhu's Sn-Log-Concavity Conjecture},
  pdfsubject={Equivariant log-concavity and orbit harmonics},
  pdfkeywords={equivariant log-concavity, orbit harmonics, set partitions, sign representation}
}

\newtheorem{theorem}{Theorem}
\newtheorem{lemma}[theorem]{Lemma}

\newcommand{\C}{\mathbb C}
\newcommand{\signrep}{\mathrm{sgn}}
\newcommand{\smult}[1]{[\signrep:#1]}
\DeclareMathOperator{\sgn}{sgn}
\DeclareMathOperator{\Stab}{Stab}

\title[Counterexamples to Zhu's log-concavity conjecture]
{Counterexamples to Zhu's
\texorpdfstring{$\mathfrak S_n$}{Sn}-Log-Concavity Conjecture}

\date{}

\begin{document}

\begin{abstract}
We disprove Conjecture~4.3 of Zhu on the
$\mathfrak S_n$-log-concavity of set-partition orbit harmonics.
For every odd $k\ge5$, set $n=k^2-5$ and $d=n-k$. If
$k+1\le m\le n$, then the sign representation has multiplicity $4$ in
$R(\Pi_{n,m})_{d-1}\otimes R(\Pi_{n,m})_{d+1}$, while its multiplicity
in $R(\Pi_{n,m})_d^{\otimes2}$ is $1$ for $k=5$ and $0$ for $k\ge7$.
Thus the equivariant injection required by log-concavity does not exist.
The first member of the family has $n=20$.
\end{abstract}

\maketitle

For $1\le m\le n$, let $\Pi_{n,m}$ be the set of unordered set
partitions of $[n]$ whose blocks have size at most $m$, and write
\[
  R_{n,m,d}=R(\Pi_{n,m})_d,
  \qquad
  X_{n,m,d}=\{P\in\Pi_{n,m}: |P|=n-d\}.
\]
Zhu proved the $\mathfrak S_n$-module isomorphism
\begin{equation}\label{eq:permutation-model}
  R_{n,m,d}\cong \C[X_{n,m,d}]
\end{equation}
and conjectured that $R(\Pi_{n,m})$ is $\mathfrak S_n$-log-concave for
$n\ge9$ \cite[Proposition~3.29 and Conjecture~4.3]{Zhu}. In particular,
log-concavity would require $\mathfrak S_n$-module injections
\[
  R_{n,m,d-1}\otimes R_{n,m,d+1}
  \hookrightarrow R_{n,m,d}^{\otimes2}
\]
for every interior degree $d$.

\begin{theorem}\label{thm:main}
Let $k\ge5$ be odd, set $n=k^2-5$ and $d=n-k$, and suppose that
$k+1\le m\le n$. Then
\[
  \smult{R_{n,m,d-1}\otimes R_{n,m,d+1}}=4,
\]
whereas
\[
  \smult{R_{n,m,d}^{\otimes2}}
  =
  \begin{cases}
    1,&k=5,\\
    0,&k\ge7.
  \end{cases}
\]
Consequently, $R(\Pi_{n,m})$ is not $\mathfrak S_n$-log-concave, and
Conjecture~4.3 of Zhu is false.
\end{theorem}

The case $k=5$ gives $n=20$, $d=15$, and every $6\le m\le20$.
No minimality in $n$ is claimed.

\begin{lemma}\label{lem:orbit-sign}
For a finite $\mathfrak S_n$-set $Y$,
\[
  \smult{\C[Y]}
  =\#\{\mathcal O\in\mathfrak S_n\backslash Y:
  \Stab(y)\subseteq\mathfrak A_n\text{ for }y\in\mathcal O\}.
\]
\end{lemma}

\begin{proof}
For an orbit $\mathcal O\cong\mathfrak S_n/H$, Frobenius reciprocity
identifies
\[
  \operatorname{Hom}_{\mathfrak S_n}(\signrep,\C[\mathcal O])
  \cong
  \operatorname{Hom}_{H}(\signrep|_H,\mathbf1).
\]
This space has dimension one if $H\subseteq\mathfrak A_n$ and zero
otherwise. Summing over the orbits proves the formula.
\end{proof}

For an automorphism $(\sigma,\pi)$ of a bipartite graph $H$, write
$\sgn_{E(H)}(\sigma,\pi)$ for the sign of its induced permutation of
$E(H)$. Let $H^\circ$ be obtained from $H$ by deleting its isolated
vertices. All graph isomorphisms below preserve the two vertex classes,
and $P_j$ denotes the path on $j$ vertices.

\begin{lemma}\label{lem:small-graphs}
The following classifications hold.
\begin{enumerate}
\item If $Z$ has four edges, then every automorphism of $Z$ acts evenly
on $E(Z)$ if and only if $Z^\circ$ is one of
\[
  C_4,\qquad P_5^{(3,2)},\qquad P_5^{(2,3)},\qquad P_4\sqcup K_2,
\]
where the superscript records the two bipartition sizes.
\item If $Z\subseteq K_{5,5}$ has five edges, then
\begin{equation}\tag{$\ast$}\label{eq:five-edge-condition}
  \sgn_{E(Z)}(\sigma,\pi)=\sgn(\sigma)\sgn(\pi)
\end{equation}
for every automorphism $(\sigma,\pi)$ of $Z$ if and only if
\[
  Z^\circ=K_{1,2}\sqcup K_{2,1}\sqcup K_2.
\]
\end{enumerate}
\end{lemma}

\begin{proof}
For the first assertion, a connected four-edge bipartite graph is
$C_4$ or a tree on five vertices. The two nonpath trees have two leaves
with a common neighbor; transposing those leaves induces one edge
transposition. The reflection of $P_5$ exchanges two pairs of edges,
and the side-preserving automorphism group of $C_4$ is generated by
row and column swaps, each of which exchanges two pairs of edges.

If $Z$ is disconnected, its component edge counts are $3+1$, $2+2$,
$2+1+1$, or $1+1+1+1$. In the first case the three-edge component is
$P_4$ or $K_{1,3}$, and only $P_4\sqcup K_2$ survives. Every other case
contains a $P_3$ or $K_{1,3}$ with an odd leaf transposition, or has two
$K_2$ components whose interchange is odd on the edges.

For the second assertion, two isolated vertices on the same side may
be transposed without moving an edge, contradicting
\eqref{eq:five-edge-condition}. Hence $Z$ uses at least four vertices
on each side. If it uses $(5,5)$ vertices, then $Z=5K_2$. If it uses
$(5,4)$ or $(4,5)$ vertices, then, up to reversing the sides,
$Z=K_{1,2}\sqcup3K_2$. Both graphs fail
\eqref{eq:five-edge-condition} by interchanging two $K_2$ components.

It remains to consider used bipartition sizes $(4,4)$. Each side then
has degree sequence $(2,1,1,1)$. Let $u$ and $v$ be the degree-two
vertices. If $uv$ is an edge, then $Z=P_4\sqcup2K_2$, which again fails
by interchanging the two $K_2$ components. If $uv$ is not an edge,
then $Z=K_{1,2}\sqcup K_{2,1}\sqcup K_2$. Its automorphism group is
generated by the two leaf transpositions. Each generator induces one
edge transposition and one transposition in exactly one vertex class,
so \eqref{eq:five-edge-condition} holds.
\end{proof}

\begin{proof}[Proof of Theorem~\ref{thm:main}]
By \eqref{eq:permutation-model},
\[
  R_{n,m,a}\otimes R_{n,m,b}
  \cong \C[X_{n,m,a}\times X_{n,m,b}].
\]
Consider $(P,Q)\in X_{n,m,a}\times X_{n,m,b}$. If blocks
$B\in P$ and $C\in Q$ satisfy $|B\cap C|\ge2$, then a transposition in
$B\cap C$ is an odd element of the stabilizer, so the orbit contributes
nothing by Lemma~\ref{lem:orbit-sign}.

Suppose instead that every intersection has size at most one. The pair
$(P,Q)$ is then encoded by a simple bipartite incidence graph $G$: its
left and right vertices are the blocks of $P$ and $Q$, and its $n$
edges are the elements of $[n]$. The $\mathfrak S_n$-orbits are the
side-preserving isomorphism classes of these graphs, and the stabilizer
is their side-preserving automorphism group acting on $E(G)$. Such an
incidence graph has no isolated vertices.

Let $G$ have $r$ left vertices and $c$ right vertices, and let
$Z=K_{r,c}\setminus G$. For an automorphism $(\sigma,\pi)$, the action
on all $rc$ cells gives
\begin{equation}\label{eq:grid-sign}
  \sgn_{E(G)}(\sigma,\pi)\,
  \sgn_{E(Z)}(\sigma,\pi)
  =\sgn(\sigma)^c\sgn(\pi)^r.
\end{equation}

For the source, $(a,b)=(d-1,d+1)$, so
\[
  (r,c)=(k+1,k-1),
  \qquad |E(Z)|=(k+1)(k-1)-n=4.
\]
Both $r$ and $c$ are even. Hence the right-hand side of
\eqref{eq:grid-sign} is $1$, and an orbit contributes precisely when
every automorphism acts evenly on $E(Z)$. Lemma~\ref{lem:small-graphs}
gives exactly four such isomorphism classes. Each listed complement
has maximum degree two, so its complement $G$ has no isolated vertices.
Moreover, every degree of $G$ is at most $k+1\le m$. Therefore
\[
  \smult{R_{n,m,d-1}\otimes R_{n,m,d+1}}=4.
\]

For the target, $(a,b)=(d,d)$, so $r=c=k$ and $|E(Z)|=5$. If $k=5$,
then \eqref{eq:grid-sign} shows that an orbit contributes exactly when
\eqref{eq:five-edge-condition} holds. By
Lemma~\ref{lem:small-graphs}, there is one such graph. Its maximum
degree is two, so its complement is an admissible incidence graph; hence
\[
  \smult{R_{n,m,d}^{\otimes2}}=1.
\]

Now let $k\ge7$. A five-edge graph uses at most five vertices on each
side, so $Z$ has two isolated vertices on each side. Transpose two
isolated left vertices and fix the right side. This automorphism acts
trivially on $E(Z)$, while \eqref{eq:grid-sign} gives
\[
  \sgn_{E(G)}=(-1)^k=-1.
\]
Thus every transverse orbit has an odd stabilizer element. Together
with the nontransverse case, Lemma~\ref{lem:orbit-sign} yields
\[
  \smult{R_{n,m,d}^{\otimes2}}=0.
\]

Finally, complex representations of $\mathfrak S_n$ are semisimple.
An injection from the source to the target would force every
irreducible multiplicity in the source to be at most the corresponding
multiplicity in the target, contrary to the sign multiplicities above.
Since the source multiplicity is nonzero, both adjacent graded pieces
are nonzero, so $d$ is an interior degree. This proves the theorem.
\end{proof}

\begin{thebibliography}{1}

\bibitem{Zhu}
H.~Zhu,
\emph{Plethysm and orbit harmonics},
\href{https://arxiv.org/abs/2602.12623v1}{arXiv:2602.12623v1 [math.CO]},
2026.

\end{thebibliography}

\end{document}